%2multibyte Version: 5.50.0.2960 CodePage: 65001
% MQF_Manuale_Master_SW55_Memoir.tex
% Master del manuale MQF costruito sullo shell SW5.5 "Memoir Class for Configurable Typesetting".
% Versione modulare: i capitoli scritti sono inclusi con \input.%
% POSIZIONE DEI FILE
% Questo master deve stare nella cartella 01_Manuale.
% I capitoli devono stare nella sottocartella 01_Manuale/Capitoli.
% Non mettere il master dentro Capitoli.
% I capitoli futuri restano come traccia, ma non generano pagine nel documento.
% Evita pagine bianche dovute all'apertura dei capitoli su pagina dispari.
% ------------------------------------------------------------
% Nomi italiani
% ------------------------------------------------------------
% ------------------------------------------------------------
% Macro matematiche ufficiali del progetto MQF
% Coerenti con MQF_Notazione.tex
% ------------------------------------------------------------
% ------------------------------------------------------------
% Ambienti di tipo teorema.
% I nomi interni restano compatibili con lo shell SW.
% Le etichette stampate sono in italiano.
% ------------------------------------------------------------

%TCIDATA{OutputFilter=latex2.dll}
%TCIDATA{Version=5.50.0.2960}
%TCIDATA{LaTeXparent=1,1,MQF_Manuale_Master.tex}


\chapter{Variabili casuali e distribuzioni}

\section{Obiettivi della lezione}

\label{sec:cap2-obiettivi}

Nel Capitolo 1 sono stati introdotti gli elementi fondamentali del linguaggio
probabilistico: lo spazio degli esiti, gli eventi, le operazioni sugli eventi
e la misura di probabilit\`{a}. Tale linguaggio consente di descrivere
l'incertezza in termini di eventi del tipo ``un certo scenario si verifica''
oppure ``un certo insieme di esiti si realizza''.

In molte applicazioni della finanza quantitativa, tuttavia, l'interesse
principale non \`{e} rivolto direttamente agli esiti elementari dello spazio
campionario, ma a quantit\`{a} numeriche che dipendono da tali esiti. Un
rendimento futuro, una perdita di portafoglio, il prezzo futuro di
un'attivit\`{a} finanziaria o il payoff di un contratto derivato sono esempi
di grandezze aleatorie. Formalmente, esse sono rappresentate mediante
variabili casuali, cio\`{e} funzioni del tipo
\[
X:\Omega\longrightarrow\mathbb{R}.
\]


Lo scopo di questo capitolo \`{e} introdurre il linguaggio delle variabili
casuali e delle loro distribuzioni. Il passaggio concettuale fondamentale
\`{e} il seguente: da una probabilit\`{a} definita sugli eventi dello spazio
campionario si costruisce una descrizione probabilistica dei valori assunti da
una grandezza numerica aleatoria. In questo modo diventa possibile studiare
eventi finanziariamente rilevanti quali
\[
\{R<0\},\qquad\{L>\ell\},\qquad\{S_{T}>K\},
\]
dove $R$ pu\`{o} rappresentare un rendimento, $L$ una perdita e $S_{T}$ il
prezzo futuro di un'attivit\`{a} finanziaria.

Al termine della lezione lo studente dovr\`{a} essere in grado di:

\begin{enumerate}
\item definire una variabile casuale come funzione misurabile dallo spazio
campionario alla retta reale;

\item interpretare eventi del tipo
\[
\{X\leq x\},\qquad\{X=x_{i}\},\qquad\{a<X\leq b\}
\]
come eventi appartenenti alla sigma-algebra $\mathcal{F}$;

\item definire la distribuzione di una variabile casuale e la funzione di
ripartizione
\[
F_{X}(x)=\mathbb{P}(X\leq x);
\]


\item distinguere tra variabili casuali discrete, descritte mediante una
funzione di massa $p_{X}$, e variabili casuali assolutamente continue,
descritte mediante una densit\`{a} $f_{X}$;

\item calcolare probabilit\`{a} associate a intervalli e code distributive,
usando la funzione di ripartizione, la funzione di massa o la densit\`{a};

\item definire e interpretare il valore atteso
\[
\mathbb{E}[X],
\]
la varianza
\[
\operatorname{Var}(X),
\]
la deviazione standard e, pi\`{u} in generale, i momenti di una variabile casuale;

\item calcolare il valore atteso di funzioni di variabili casuali, con
particolare riferimento a payoff finanziari del tipo $g(X)$;

\item introdurre il concetto di quantile, con particolare attenzione alla sua
interpretazione nella misura delle perdite e del rischio;

\item riconoscere alcune famiglie distributive fondamentali, tra cui le
distribuzioni di Bernoulli, binomiale, di Poisson, uniforme, normale e lognormale.
\end{enumerate}

Il capitolo costruisce quindi il ponte tra la teoria degli eventi e la
misurazione quantitativa del rischio. La sequenza logica che guider\`{a}
l'esposizione \`{e}
\[
(\Omega,\mathcal{F},\mathbb{P}) \quad\longrightarrow\quad X \quad
\longrightarrow\quad F_{X},\ p_{X},\ f_{X} \quad\longrightarrow\quad
\mathbb{E}[X],\ \operatorname{Var}(X),\ q_{\alpha}(X).
\]
Questa sequenza permette di passare dalla descrizione astratta dell'incertezza
alla rappresentazione analitica di rendimenti, perdite e payoff aleatori.

\section{Motivazione finanziaria: rendimenti, perdite e payoff aleatori}

\label{sec:cap2-motivazione-finanziaria}

Nelle applicazioni finanziarie l'incertezza non riguarda soltanto il
verificarsi o meno di un evento, ma soprattutto il valore numerico assunto da
una grandezza economica o finanziaria. L'esito futuro di un mercato pu\`{o}
determinare un rendimento positivo o negativo, una perdita di portafoglio, il
prezzo futuro di un'attivit\`{a} finanziaria o il payoff di un contratto
derivato. Tali grandezze non sono note al tempo iniziale, ma diventano note
solo dopo la realizzazione dello scenario rilevante.

Il linguaggio degli eventi, introdotto nel Capitolo 1, consente di formulare
affermazioni del tipo
\[
A=\{\text{il mercato scende}\}, \qquad B=\{\text{il debitore non rimborsa il
prestito}\}.
\]
Tuttavia, in molti problemi di finanza quantitativa, \`{e} necessario
descrivere non solo se un evento si verifica, ma anche quanto vale una
determinata grandezza in ciascuno scenario. Per questo motivo si introduce il
concetto di variabile casuale.

Ad esempio, se $R$ indica il rendimento aleatorio di un'attivit\`{a}
finanziaria, non interessa soltanto l'evento
\[
\{R<0\},
\]
cio\`{e} l'evento in cui il rendimento \`{e} negativo. Interessa anche la
distribuzione dei possibili valori di $R$, perch\'{e} tale distribuzione
consente di valutare la probabilit\`{a} di perdite moderate, di perdite
estreme, di rendimenti elevati o di risultati prossimi al valore medio.

In modo analogo, se $L$ indica la perdita aleatoria di un portafoglio, eventi
del tipo
\[
\{L>\ell\}
\]
descrivono il superamento di una soglia di perdita $\ell$. La quantit\`{a}
\[
\mathbb{P}(L>\ell)
\]
rappresenta quindi una probabilit\`{a} di coda, cio\`{e} la probabilit\`{a}
che la perdita superi un livello prefissato. Questo tipo di probabilit\`{a}
\`{e} centrale nella misurazione del rischio finanziario, poich\'{e} permette
di studiare non soltanto il comportamento medio di un portafoglio, ma anche la
severit\`{a} degli scenari sfavorevoli.

Un terzo esempio fondamentale \`{e} fornito dai contratti derivati. Sia
$S_{T}$ il prezzo aleatorio, al tempo futuro $T$, di un'attivit\`{a}
sottostante. Il payoff di un'opzione call europea con prezzo di esercizio $K$
\`{e}
\[
H=(S_{T}-K)^{+},
\]
dove
\[
(S_{T}-K)^{+}=\max\{S_{T}-K,0\}.
\]
Anche in questo caso il payoff $H$ \`{e} una variabile casuale: il suo valore
dipende dallo scenario futuro e, in particolare, dal valore assunto da $S_{T}%
$. L'evento
\[
\{S_{T}>K\}
\]
indica che l'opzione termina in the money, mentre la quantit\`{a}
\[
\mathbb{P}(S_{T}>K)
\]
misura la probabilit\`{a} che il payoff sia strettamente positivo.

Questi esempi mostrano che una parte rilevante della finanza quantitativa
pu\`{o} essere formulata mediante grandezze aleatorie numeriche:
\[%
\begin{array}
[c]{ll}%
R=\mbox{rendimento aleatorio}, & L=\mbox{perdita aleatoria},\\[4pt]%
S_{T}=\mbox{prezzo futuro aleatorio}, & H=\mbox{payoff aleatorio}.
\end{array}
\]
La variabile casuale \`{e} dunque lo strumento matematico che consente di
associare un numero reale a ciascuno scenario possibile.

Il problema centrale diventa allora descrivere la distribuzione probabilistica
di tali grandezze. In particolare, si vogliono rispondere domande del tipo:
\[
\mathbb{P}(R<0)=?
\]
\[
\mathbb{P}(L>\ell)=?
\]
\[
\mathbb{E}[R]=?
\]
\[
\operatorname{Var}(R)=?
\]
\[
q_{\alpha}(L)=?
\]
Le prime due domande riguardano probabilit\`{a} di eventi definiti tramite
variabili casuali. La terza riguarda il rendimento medio atteso. La quarta
riguarda la dispersione del rendimento intorno al proprio valore atteso. La
quinta riguarda un quantile della distribuzione della perdita, cio\`{e} una
soglia che viene superata solo con una probabilit\`{a} prefissata.

In questa prospettiva, una variabile casuale non \`{e} semplicemente un numero
incerto. Essa \`{e} una funzione che trasforma gli esiti elementari in valori
numerici:
\[
X:\Omega\longrightarrow\mathbb{R}.
\]
La probabilit\`{a}, definita originariamente sugli eventi dello spazio
campionario, induce cos\`{\i} una distribuzione sui valori assunti da $X$.
Questo passaggio \`{e} essenziale: permette di sostituire la descrizione
astratta degli scenari con una descrizione quantitativa dei risultati
economici associati a tali scenari.

Il capitolo sviluppa quindi gli strumenti necessari per descrivere e
utilizzare questa distribuzione: la funzione di ripartizione $F_{X}$, la
funzione di massa $p_{X}$ nel caso discreto, la densit\`{a} $f_{X}$ nel caso
continuo, il valore atteso, la varianza, i momenti e i quantili. Questi
oggetti costituiscono la base probabilistica per l'analisi di rendimenti,
perdite, payoff e rischio finanziario.

\section{Dallo spazio degli eventi alle variabili casuali}

\label{sec:cap2-dallo-spazio-eventi-vc}

Nel Capitolo 1 l'incertezza \`{e} stata descritta mediante uno spazio di
probabilit\`{a}
\[
(\Omega,\mathcal{F},\mathbb{P}),
\]
dove $\Omega$ rappresenta l'insieme degli esiti possibili, $\mathcal{F}$ \`{e}
la sigma-algebra degli eventi osservabili e $\mathbb{P}$ \`{e} la misura di
probabilit\`{a} definita su tali eventi.

In molte applicazioni, tuttavia, non interessa direttamente l'esito elementare
$\omega\in\Omega$, ma una quantit\`{a} numerica che dipende da tale esito. Una
variabile casuale formalizza precisamente questo passaggio.

\begin{definition}
[Variabile casuale]Sia $(\Omega,\mathcal{F},\mathbb{P})$ uno spazio di
probabilit\`{a}. Una variabile casuale reale \`{e} una funzione
\[
X:\Omega\longrightarrow\mathbb{R}
\]
tale che, per ogni $x\in\mathbb{R}$,
\[
\{\omega\in\Omega:X(\omega)\leq x\}\in\mathcal{F}.
\]

\end{definition}

La condizione precedente significa che, per ogni soglia reale $x$, l'insieme
degli esiti per i quali la variabile casuale assume un valore non superiore a
$x$ \`{e} un evento. Di conseguenza, la sua probabilit\`{a} \`{e} ben definita.

Per alleggerire la notazione, si scrive normalmente
\[
\{X\leq x\}
\]
al posto di
\[
\{\omega\in\Omega:X(\omega)\leq x\}.
\]
Con questa convenzione, la quantit\`{a}
\[
\mathbb{P}(X\leq x)
\]
deve essere letta come abbreviazione di
\[
\mathbb{P}(\{\omega\in\Omega:X(\omega)\leq x\}).
\]


In modo analogo, scritture come
\[
\{X=x_{i}\},\qquad\{a<X\leq b\},\qquad\{X\in B\}
\]
indicano rispettivamente gli eventi
\[
\{\omega\in\Omega:X(\omega)=x_{i}\},
\]
\[
\{\omega\in\Omega:a<X(\omega)\leq b\},
\]
\[
\{\omega\in\Omega:X(\omega)\in B\},
\]
dove $B$ \`{e} un sottoinsieme opportuno della retta reale.

\paragraph{Esempio guida: rendimento discreto su tre regimi di mercato.}

Consideriamo un modello elementare con sei esiti possibili:
\[
\Omega=\{\omega_{1},\omega_{2},\omega_{3},\omega_{4},\omega_{5},\omega_{6}\}.
\]
Gli esiti elementari rappresentano sei scenari di mercato. Supponiamo tuttavia
che l'informazione rilevante distingua soltanto tre regimi: mercato
sfavorevole, mercato intermedio e mercato favorevole. Introduciamo quindi la
partizione
\[
A_{1}=\{\omega_{1},\omega_{2}\}, \qquad A_{2}=\{\omega_{3},\omega_{4}\},
\qquad A_{3}=\{\omega_{5},\omega_{6}\}.
\]
L'evento $A_{1}$ rappresenta il regime sfavorevole, $A_{2}$ il regime
intermedio e $A_{3}$ il regime favorevole. La sigma-algebra generata da questa
partizione \`{e}
\[
\mathcal{F}=\sigma(A_{1},A_{2},A_{3}).
\]
Esplicitamente,
\[
\mathcal{F}= \{\varnothing, A_{1},A_{2},A_{3}, A_{1}\cup A_{2},A_{1}\cup
A_{3},A_{2}\cup A_{3}, \Omega\}.
\]
In questo modello non tutte le distinzioni tra gli esiti elementari sono
osservabili attraverso $\mathcal{F}$: gli eventi ammissibili sono le unioni
degli elementi della partizione.

Definiamo ora una variabile casuale $X$ che rappresenta il rendimento di un
portafoglio in ciascuno scenario. Supponiamo che il rendimento dipenda
soltanto dal regime di mercato e sia dato dalla seguente tabella:
\[%
\begin{array}
[c]{c|c|c}%
\mbox{esito} & \mbox{regime} & X(\omega)\\\hline
\omega_{1},\omega_{2} & A_{1}\ \mbox{sfavorevole} & -0.04\\
\omega_{3},\omega_{4} & A_{2}\ \mbox{intermedio} & 0.01\\
\omega_{5},\omega_{6} & A_{3}\ \mbox{favorevole} & 0.06
\end{array}
\]
La variabile casuale $X$ trasforma quindi i sei esiti elementari in tre valori
reali:
\[
x_{1}=-0.04, \qquad x_{2}=0.01, \qquad x_{3}=0.06.
\]
I tre valori assunti da $X$ determinano una partizione sulla retta reale. Le
regioni rilevanti sono
\[
(-\infty,-0.04), \qquad[-0.04,0.01), \qquad[0.01,0.06, \qquad[0.06,+\infty).
\]
Eventi descritti in termini di rendimenti (valori numerici) $X \in\mathbb{R}$
sono in corrispondenza con eventi qualitativi in $\Omega$. Ad esempio,
l'evento ``rendimento negativo'' \`{e}
\[
\{X<0\}=A_{1}=\{\omega_{1},\omega_{2}\}.
\]
L'evento ``rendimento almeno pari all'$1\%$'' \`{e}
\[
\{X\geq0.01\}=A_{2}\cup A_{3} =\{\omega_{3},\omega_{4},\omega_{5},\omega
_{6}\}.
\]
L'evento ``rendimento non superiore all'$1\%$'' \`{e}
\[
\{X\leq0.01\}=A_{1}\cup A_{2} =\{\omega_{1},\omega_{2},\omega_{3},\omega
_{4}\}.
\]
In tutti questi casi la probabilit\`{a} resta formalmente una probabilit\`{a}
assegnata a eventi appartenenti a $\mathcal{F}$:
\[
\mathbb{P}(X<0)=\mathbb{P}(A_{1}),
\]
\[
\mathbb{P}(X\geq0.01)=\mathbb{P}(A_{2}\cup A_{3}),
\]
\[
\mathbb{P}(X\leq0.01)=\mathbb{P}(A_{1}\cup A_{2}).
\]
La variabile casuale consente per\`{o} di descrivere tali eventi mediante
condizioni sui valori numerici del rendimento.

\paragraph{Controimmagini e misurabilit\`{a}.}

Il legame formale tra valori numerici ed eventi \`{e} dato dalla nozione di
controimmagine. Sia
\[
X:\Omega\longrightarrow\mathbb{R}
\]
una funzione e sia $B\subseteq\mathbb{R}$. Si definisce controimmagine di $B$
mediante $X$ l'insieme
\[
X^{-1}(B)=\{\omega\in\Omega:X(\omega)\in B\}.
\]
La controimmagine $X^{-1}(B)$ \`{e} quindi un sottoinsieme di $\Omega$: essa
raccoglie tutti gli esiti elementari che, attraverso la funzione $X$, vengono
mandati in $B$.

La notazione $X^{-1}(B)$ indica l'insieme degli esiti la cui immagine
appartiene a $B$. Ad esempio, se $B=(-\infty,a]$, allora
\[
X^{-1}((-\infty,a]) = \{\omega\in\Omega:X(\omega)\leq a\} = \{X\leq a\}.
\]
Analogamente,
\[
X^{-1}((a,b]) = \{\omega\in\Omega:a<X(\omega)\leq b\} = \{a<X\leq b\}.
\]


La nozione di variabile casuale richiede che queste controimmagini siano
eventi ai quali sia possibile assegnare una probabilit\`{a}. Per garantire
questo fatto si dimostra che \`{e} necessario e sufficiente che, per ogni
scelta arbitraria del valore $a\in\mathbf{R}$ negli intervalli di tipo
$(-\infty,a],$ risulti:%
\begin{equation}
X^{-1}(-\infty,a]\in\mathcal{F}.\label{cmisur}%
\end{equation}
Se vale la precedente condizione, allora discende anche
\[
X^{-1}(a,b]\in\mathcal{F}%
\]
per ogni scelta di $a,b$.  In termini generali, una funzione $X:\Omega
\rightarrow\mathbb{R}$ per cui vale la condizione \ref{cmisur} \`{e} detta
variabile casuale \textbf{misurabile}.

Nel seguito useremo prevalentemente eventi del tipo
\[
\{X\leq x\},\qquad\{a<X\leq b\},\qquad\{X\in B\},
\]
poich\'{e} sono quelli che permettono di costruire la funzione di
ripartizione, la funzione di massa, la densit\`{a} e i quantili.

\begin{example}
[Controimmagini nell'esempio del rendimento discreto]Riprendiamo la variabile
casuale $X$ definita sull'insieme
\[
\Omega=\{\omega_{1},\omega_{2},\omega_{3},\omega_{4},\omega_{5},\omega_{6}\}
\]
mediante la tabella
\[%
\begin{array}
[c]{c|c|c}%
\mbox{esito} & \mbox{regime} & X(\omega)\\\hline
\omega_{1},\omega_{2} & A_{1}\ \mbox{sfavorevole} & -0.04\\
\omega_{3},\omega_{4} & A_{2}\ \mbox{intermedio} & 0.01\\
\omega_{5},\omega_{6} & A_{3}\ \mbox{favorevole} & 0.06
\end{array}
\]
dove
\[
A_{1}=\{\omega_{1},\omega_{2}\}, \qquad A_{2}=\{\omega_{3},\omega_{4}\},
\qquad A_{3}=\{\omega_{5},\omega_{6}\}.
\]
La sigma-algebra considerata \`{e}
\[
\mathcal{F}=\sigma(A_{1},A_{2},A_{3}) = \{\varnothing, A_{1},A_{2},A_{3},
A_{1}\cup A_{2},A_{1}\cup A_{3},A_{2}\cup A_{3}, \Omega\}.
\]


Poich\'{e} i possibili rendimenti $X$ assumo soltanto i valori%
\[
-0.04<0.01<0.06,
\]
la controimmagine di un intervallo del tipo $(-\infty,a]$ dipende dalla
posizione della soglia $a$ rispetto a questi tre valori. Si ottiene:
\[%
\begin{array}
[c]{c|c}%
\text{valori }a\text{ per intervalli }(-\infty,a] & X^{-1}((-\infty
,a])=\{X\leq a\}\\\hline
a<-0.04 & \varnothing\\
-0.04\leq a<0.01 & A_{1}\\
0.01\leq a<0.06 & A_{1}\cup A_{2}\\
a\geq0.06 & \Omega
\end{array}
\]
Ogni insieme nella seconda colonna appartiene a $\mathcal{F}$. La tabella
mostra quindi, in modo esplicito, che $X$ \`{e} misurabile rispetto alla
sigma-algebra considerata.

La stessa logica vale per altri sottoinsiemi della retta reale. Ad esempio,
\[
X^{-1}([0,+\infty))=A_{2}\cup A_{3},
\]
perch\'{e} i valori non negativi di $X$ sono $0.01$ e $0.06$. Inoltre,
\[
X^{-1}((-\infty,0))=A_{1}.
\]
In termini finanziari, il primo evento rappresenta gli scenari con rendimento
non negativo, mentre il secondo rappresenta gli scenari con rendimento negativo.
\end{example}

\paragraph{Variabili indicatrici.}

Un caso particolarmente importante di variabile casuale \`{e} la variabile
indicatrice di un evento. Se $A\in\mathcal{F}$, si definisce
\[
\mathbf{1}_{A}(\omega)=
\begin{cases}
1, & \mbox{se }\omega\in A,\\
0, & \mbox{se }\omega\notin A.
\end{cases}
\]
La variabile $\mathbf{1}_{A}$ assume quindi valore $1$ quando l'evento $A$ si
verifica e valore $0$ quando l'evento non si verifica.

Nell'esempio del rendimento discreto, l'indicatrice dell'evento di rendimento
negativo \`{e}
\[
\mathbf{1}_{\{X<0\}}=\mathbf{1}_{A_{1}}.
\]
Essa vale $1$ negli scenari $\omega_{1}$ e $\omega_{2}$, e vale $0$ negli
altri scenari. Analogamente, l'indicatrice dell'evento di rendimento almeno
pari all'$1\%$ \`{e}
\[
\mathbf{1}_{\{X\geq0.01\}}=\mathbf{1}_{A_{2}\cup A_{3}}.
\]


Nelle applicazioni finanziarie, variabili indicatrici di questo tipo compaiono
frequentemente:
\[
\mathbf{1}_{\{R<0\}}, \qquad\mathbf{1}_{\{L>\ell\}}, \qquad\mathbf{1}%
_{\{S_{T}>K\}}.
\]
Esse indicano rispettivamente il verificarsi di un rendimento negativo, il
superamento di una soglia di perdita e il superamento dello strike da parte
del sottostante. Saranno utili anche per rappresentare payoff digitali, eventi
di default e vincoli probabilistici.

In conclusione, una variabile casuale non sostituisce lo spazio di
probabilit\`{a}, ma lo utilizza per costruire una descrizione numerica
dell'incertezza. Lo spazio $(\Omega,\mathcal{F},\mathbb{P})$ resta il
fondamento formale; la variabile casuale $X$ traduce gli esiti $\omega$ in
valori reali $X(\omega)$; la distribuzione di $X$, introdotta nella sezione
successiva, descrive la probabilit\`{a} con cui tali valori possono manifestarsi.